\section{Introduction}
\label{sec:introduction}

Large language models are often evaluated on whether they can produce an answer. In deployment, that objective can be wrong. For research assistants, search tools, tutoring systems, and agent pipelines, a confident but fabricated answer can be more harmful than an explicit ``I don't know.'' Our main question is simple: when a model is uncertain, or when a request is impossible to satisfy, does prompt pressure to answer push it toward compliant-looking output instead of honest abstention?

This question matters because many real systems implicitly reward visible action. A downstream component may treat any non-empty response as progress, even if the response should have been an abstention. If models learn to satisfy the surface form of a request while ignoring the underlying epistemic constraint, then the failure is not just hallucination in the narrow sense. It is a broader answer-pressure failure in which the model exploits an easy loophole: answering anyway.

Existing work provides many pieces of this picture, but not a unified test. Research on strategic deception shows that models can preserve undesirable behavior under training pressure \citep{hubinger2024sleeper}. Truthfulness and factuality benchmarks measure plausible falsehoods and calibrated answering \citep{lin2021truthfulqa,wei2024simpleqa}. Hallucination benchmarks study fabricated content and verifier design \citep{li2023halueval,hu2024refchecker,ravichander2025halogen}. AbstentionBench shows that frontier reasoning models still fail badly on unanswerable questions \citep{kirichenko2025abstentionbench}. However, these lines of work usually evaluate one failure mode at a time. They do not directly ask whether the same answer-pressure dynamic connects factual guessing, abstention failure, and fabricated compliance on impossible requests.

We address that gap with a controlled cross-benchmark study. We hold question sets fixed and vary only the instruction pressure around answering. We compare a neutral \baseline prompt, an \abstainallowed prompt that explicitly permits ``I don't know,'' and a \forcedanswer prompt that forbids abstention. We test \gptfourone and \gptfivemini on \simpleqa, \truthfulqa, and false-presupposition prompts from \halogen, where the impossible subset gives a clean setting in which abstention is the correct behavior.

The main result is sharp. On impossible \halogen prompts, \forcedanswer raises the loophole-attempt rate from 28.3\% to 100\% for \gptfourone and from 31.7\% to 100\% for \gptfivemini. Fabricated attempts rise by 75.0 and 61.7 percentage points, respectively. On factual benchmarks, the pattern depends strongly on the model. On \simpleqa, \gptfivemini moves from an 8\% attempt rate under \abstainallowed to 100\% under \forcedanswer, with an 82-point increase in incorrect answers. \gptfourone changes little because it already tends to answer even when abstention is encouraged.

Our contributions are:
\begin{itemize}
    \item We define \loopholeuse as attempted compliance when abstention is epistemically correct, which lets us compare impossible requests and factual uncertainty within one framework.
    \item We conduct a controlled study across three benchmark families and three prompt conditions, totaling 960 model generations and 600 grader calls.
    \item We show that impossible structured prompts produce the strongest loophole effect, with both tested models switching to universal attempted compliance under \forcedanswer.
    \item We identify a strong model difference on factual QA: \gptfivemini exhibits a large abstention-versus-bluffing tradeoff, while \gptfourone mostly answers regardless of instruction.
\end{itemize}

The rest of the paper is organized as follows. \Secref{sec:related_work} situates the study in prior work on deception, truthfulness, hallucination, and abstention. \Secref{sec:methodology} describes the benchmark slices, prompt conditions, grading pipeline, and statistical tests. \Secref{sec:results} reports the main quantitative findings, and \Secref{sec:discussion} discusses interpretation, limitations, and broader implications.
